\section{Results}
\label{sec:results}

Every result in this section is for the simulated scenario of Sec.~\ref{sec:system} (synthetic road grid, 20 RSUs, 60 vehicles, injected degradation, 5 seeds); we repeat the scope in each subsection.
Differences are absolute differences in post-onset deadline success (steps $\ge 400$) unless stated, written as mean [95\% interval over 5 seeds], with the label of the pre-registered rule where one applies (Sec.~\ref{sec:protocol}).
The shared operating point is $\rho=1.0$ with 20\% of RSUs degraded; it belongs to both sweeps.

\subsection{Behavioural trust in closed-loop dispatch (H1, H3)}
\label{sec:behav}

\paragraph*{Task success against no trust}
In this simulation, both trust heads raise post-onset deadline success over the tuned no-trust baseline A at all four degraded fractions (5\%, 10\%, 20\%, 30\%; $\rho=1.0$), with non-overlapping seed error bars; at every point the paired interval over the 5 seeds lies above $0$ with a mean above the $0.01$ floor, so the pre-registered label is WIN at 4 of 4 points for each head, and H3 holds against A.
% src: A5; F26 (table and paired effects), F27 (verdicts); D56
The gains over A are $+0.021$, $+0.038$, $+0.050$, $+0.054$ for C and $+0.016$, $+0.029$, $+0.034$, $+0.039$ for D at 5\%, 10\%, 20\%, 30\% (Table~\ref{tab:paper_fraction_effects}).
% src: A5 (C: +0.021, +0.038, +0.050, +0.054; D: +0.016, +0.029, +0.034, +0.039); results/s10/tables/paper_fraction_effects
At the shared point the post-onset success is $0.871\pm0.003$ for A, $0.921\pm0.010$ for C and $0.905\pm0.008$ for D (Table~\ref{tab:paper_ablation_shared_point}, Figs.~\ref{fig:frac} and~\ref{fig:ablation}).
% src: A5 (A 0.871 +/- 0.003, C 0.921 +/- 0.010, D 0.905 +/- 0.008); results/s10/tables/ablation_shared_point
B, the certificate-only variant, makes the same decisions as A in this scenario (asserted on all 2000 paired draws), so its row is A's row; the comparison is against no behavioural signal.
% src: A4, X4

\input{\tabdir/paper_fraction_effects.tex}
\input{\tabdir/paper_ablation_shared_point.tex}

\begin{figure*}[t]
\centering
\includegraphics[width=0.95\textwidth]{\figdir/fig3_success_vs_fraction.pdf}
\caption{Post-onset deadline success against the degraded fraction at $\rho=1.0$, in this simulation. (a) Mean $\pm$ sd over 5 seeds. (b) Paired gain over A with the 95\% interval over seeds and the pre-registered floor of $0.01$. B equals A; B-oracle is an upper bound; Gate-C and H are extra baselines. (Generated by \texttt{scripts/make\_figures.py}.)}
\label{fig:frac}
\end{figure*}

\begin{figure*}[t]
\centering
\includegraphics[width=0.95\textwidth]{\figdir/fig4_ablation.pdf}
\caption{Ablation at the shared point ($\rho=1.0$, 20\% degraded), in this simulation. Left: post-onset success; middle: dispatches to degraded RSUs from onset to avoidance (\nav); right: the same over the fixed window $[400,1200)$ (\nwin). C and D appear together; D did not beat C. (Generated by \texttt{scripts/make\_figures.py}.)}
\label{fig:ablation}
\end{figure*}

\paragraph*{Against the oracle}
In this simulation, C is within $0.001$--$0.003$ of B-oracle's success at every fraction; the label is NULL, because the oracle's advantage does not reach the $0.01$ floor.
D is $0.006$--$0.018$ below the oracle and is NEGATIVE against it at 10\%, 20\% and 30\%.
% src: A6; F27 (item 2), F26; Table paper_fraction_effects (C minus B-oracle -0.0012..-0.0028; D minus B-oracle -0.0061..-0.0184)
B-oracle reads the sealed ground truth, so it is an upper bound on what any behavioural signal could add, not a competing method.
% src: A6, X4

\paragraph*{Dispatches to degraded RSUs (H1)}
In this simulation, we always state the three dispatch quantities together: \nav, the dispatches to degraded RSUs from onset to the run's own effective-avoidance step; \nwin, the dispatches over the fixed window $[400,1200)$, which does not depend on the avoidance rule; and the never-avoided rate, the fraction of runs that never reach the avoidance rule.
% src: A7, R2, X9
At the shared point these are, for A, $627.5\pm100.4$ / $912\pm21$ / $0.47\pm0.12$; for C, $153\pm48$ / $506\pm75$ / $0.03\pm0.04$; for D, $210.6\pm84.5$ / $610\pm84$ / $0.03\pm0.04$; and for B-oracle, $24\pm26$ / $395\pm90$ / $0.01\pm0.03$.
% src: A7 (A n_av 627.5 +/- 100.4, n_win 912 +/- 21, never 0.47 +/- 0.12; C 153 +/- 48, 506 +/- 75, 0.03 +/- 0.04; D 210.6 +/- 84.5, 610 +/- 84, 0.03 +/- 0.04; B-oracle 24 +/- 26, 395 +/- 90, 0.01 +/- 0.03); F26; Table ablation_shared_point
Relative to A, in units of A's count, the \nav\ reduction is $+0.756$ [$+0.676$, $+0.836$] for C and $+0.664$ [$+0.523$, $+0.806$] for D (WIN for both), and the fixed-window reduction is smaller, $+0.446$ [$+0.351$, $+0.541$] and $+0.331$ [$+0.221$, $+0.442$].
% src: A7 (D57 normalisation); F27 item 3
A reaches the avoidance rule in only about half of its runs (never-avoided $0.47$), so its \nav\ is partly a full-window count; the fixed-window reduction is the more conservative statement.
% src: A7 (never-avoided 0.47); F27 item 3 (A reaches the rule in only 53% of runs)
Against B-oracle both heads are NEGATIVE, the oracle being better by $0.205$ (C) and $0.296$ (D) of A's count.
% src: A7
Detection latency is descriptive and never gates a verdict: C detects degradation $151\pm22$ steps after onset and D $220\pm25$ steps; A and B cannot detect by construction; the oracle detects at once by construction; and D leaves 4\% of degraded RSUs undetected inside the 400-step window (detected fraction $0.96\pm0.07$).
% src: A7 (C 151 +/- 22, D 220 +/- 25, detected fraction 0.96 +/- 0.07); F26, F27
Latency is compared under one shared rule, not at a detection operating point; the rule's false-alarm rate was 0.32--0.44 of healthy RSUs where measured (S16) and was not measured elsewhere.
% src: A7 (D78 shared-rule caveat, verbatim); F51

\paragraph*{The shift is partial}
Trust is one weighted term against latency and load, so degraded RSUs still receive a share of tasks.
In the demo scenario (seed 100, draw 7, a tuning seed; an illustration, not a result) the share of plateau dispatches to degraded RSUs is $0.226$ for A, $0.128$ for C and $0.142$ for D, so degraded RSUs still receive about 13--14\% of plateau tasks under C and D.
% src: A8 (F22 demo: 0.226, 0.128, 0.142); X16 (labelled illustration)
Figure~\ref{fig:timeline} shows that dispatch moves away from the degrading segment while trust falls.
% src: A8 ("dispatch moves away from the degrading segment while trust falls (Figure 2, both panels)"); no value is read off the figure (TODO if wanted)

\begin{figure*}[t]
\centering
\includegraphics[width=0.85\textwidth]{\figdir/fig2_trust_timeline.pdf}
\caption{Trust falls on the degrading segment while dispatch moves away from it (closed loop, D against A), in this simulation. (a) One scenario: the demo scenario (seed 100, draw 7, a tuning seed), shown as an illustration. (b) 5 seeds $\times$ 50 draws, mean $\pm$ sd over seeds. (Generated by \texttt{scripts/make\_figures.py}.)}
\label{fig:timeline}
\end{figure*}

\paragraph*{Along the correlation sweep}
Table~\ref{tab:paper_rho_sweep} gives post-onset success along the correlation sweep at 20\% degraded.
In this simulation, the labels of C and D against A are WIN at every point, and against the oracle C is NULL on success and D is NEGATIVE.
% src: F27 (every other point has the same labels as the shared point for vs-A and vs-oracle); results/s10/tables/paper_rho_sweep
Gate-C and H were run at the four fraction-sweep points only (D59, D65) and were not run at the correlation-sweep points.
% src: CLAIMS "items the paper draft may not fill": Gate-C and H at rho-sweep points not run

\input{\tabdir/paper_rho_sweep.tex}

\subsection{How much of the gain needs learning? (H against C)}
\label{sec:learning}

H uses each RSU's own \texttt{success\_ewma} as its trust: the same fixed rule, no gate, no model and no neighbour information.
H's weights $(0.5,0.25,0.3)$ were tuned once on seed 100 by the same procedure as every other variant and applied unchanged.
% src: A23 (D65, F32); R7
In this simulation, H came within $0.003$--$0.005$ of the learned node-local head C on post-onset success at all four degraded fractions: C minus H is $+0.0034$ [$+0.0027$, $+0.0041$] at 5\%, $+0.0035$ [$+0.0028$, $+0.0042$] at 10\%, $+0.0038$ [$+0.0032$, $+0.0045$] at 20\% (the shared point) and $+0.0045$ [$+0.0033$, $+0.0057$] at 30\%.
% src: A23 (verbatim differences and intervals); F33, F34; Table paper_fraction_effects
Every interval lies above $0$ but below the $0.01$ floor, so the pre-registered label is NULL at all four points and the verdict is NOT SEPARATED (WIN at 0 of 4, NEGATIVE at 0 of 4); a difference of $0.025$ or more in either direction is ruled out.
% src: A23 (D65 verdict NOT SEPARATED; 0.025 ruled out)
At the shared point \nav\ / \nwin\ / never-avoided are $152.7$ / $506.3$ / $0.028$ for C, $155.1$ / $510.7$ / $0.028$ for H and $627.5$ / $912.4$ / $0.468$ for A; C detects $5$--$10$ steps earlier than H with every interval including $0$ (descriptive; latency is compared under one shared rule, not at a detection operating point; the rule's false-alarm rate was 0.32--0.44 of healthy RSUs where measured (S16) and was not measured elsewhere).
% src: A23 (R2: always together; D78 caveat inside the parenthesis, as A23 words it); F33, F34 items 3-4
Every point has 250 draws attempted, 0 failed and 250 aggregated, and C recomputed bit-for-bit equal to its earlier result.
% src: A23 (D42)
Descriptively, a non-learned behavioural monitor recovers most of the trust-aware gain: H minus A is $+0.0177$, $+0.0344$, $+0.0460$, $+0.0493$ and C minus A is $+0.0211$, $+0.0379$, $+0.0498$, $+0.0538$ at 5\%, 10\%, 20\%, 30\% (Table~\ref{tab:paper_fraction_effects}).
% src: A27 (descriptive; the pre-registered comparison is A23); F33
The pre-registered comparison is the NULL above; this last sentence describes the same data and is not a second test.

\subsection{Weighted rule versus gate}
\label{sec:gate}

Gate-C removes candidates below $\tau=0.9$ from C's trust before the same rule selects among the survivors; $\tau$ was chosen on one seed from a flat grid on which all values lay within $0.0007$ of each other.
% src: A15 (tau chosen on one seed from a flat grid: all values within 0.0007), D59
In this simulation, the gate performed equivalently to the weighted rule on task success: C minus Gate-C is $+0.0003$ [$-0.0003$, $+0.0009$] at 5\%, $+0.0001$ [$-0.0004$, $+0.0005$] at 10\%, $-0.0002$ [$-0.0006$, $+0.0003$] at 20\% and $-0.0003$ [$-0.0006$, $+0.0001$] at 30\%.
The label is NULL at all four points and the pre-registered verdict is NOT SEPARATED.
% src: A15; F28-F31, D59
Gate-C sends slightly fewer tasks to degraded RSUs than C at the shared point (\nwin\ $493.5$ against $506.3$) without a success difference; this is descriptive, as D59 judged success only.
% src: A15 (n_win 493.5 vs 506.3)
The gate acted often---it removed a candidate on 31--44\% of decisions---and its empty-candidate-set fallback rate at steps $\ge 400$ was $4.8\%$, $6.2\%$, $12.1\%$ and $19.1\%$ from 5\% to 30\% degraded.
% src: A15 (31-44%, fallback rates 4.8%, 6.2%, 12.1%, 19.1%)
That trust enters as a weighted term and not as a gate is therefore a \emph{design argument}, not a measured gain: the weighted form's benefits are an exactly decomposable score (Layer 1) and no empty-candidate-set failure mode.
% src: A15 (D62 framing), X3

\subsection{Graph structure}
\label{sec:graph}

The load-bearing hypothesis was H2: that message passing beats a node-local model.
We tested it four times, each with a rule fixed in advance: under shared onset (S4), under staggered onset (S13), with a contained graph head (S14), and with that head and the node-local head both trained on a forward-looking objective (S15, Sec.~\ref{sec:objective}); the closed loop (S9) was a secondary check (Table~\ref{tab:paper_graph_tests}).
None supports it: the graph head was worse than the node-local head in the first two tests and not distinguishable from it in the last two.
% src: A9, A10, A25, A28, A29 (X1, X2: tested and not supported; X23: no claim about other objectives)

\input{\tabdir/paper_graph_tests.tex}

\paragraph*{Shared onset (S4)}
On the pre-registered diagnostic AUC, the 2-layer GNN head D did \emph{not} beat the capacity-matched node-local head C.
In this simulation, the paired difference D minus C is negative at every $\rho$: $-0.0297$ [$-0.0542$, $-0.0051$] at $\rho=1.0$ and $-0.0517$ [$-0.0902$, $-0.0131$] at $\rho=0$; the slope of the gap on realised correlation is $+0.0300$ [$-0.0017$, $+0.0617$], not distinguishable from zero.
The verdict is NEGATIVE (C beats D).
% src: A9 (F16, F17; D43, D47, D49); Figure 5
Collusion is NOISY at both $\rho$, which is inconclusive: the models never saw colluders in training.
Cold start is NULL at $\rho=1.0$ ($-0.0273$ [$-0.0935$, $+0.0389$]; an advantage of $0.05$ is ruled out), and at $\rho=0$ D is worse and the 95\% interval over seeds excludes $0$ ($-0.1431$ [$-0.2226$, $-0.0635$]; reported as information).
C detects degradation about $77$--$97$ steps earlier than D at every $\rho$.
Latency is compared under one shared rule, not at a detection operating point; the rule's false-alarm rate was 0.32--0.44 of healthy RSUs where measured (S16) and was not measured elsewhere.
% src: A9 (D78 shared-rule caveat, verbatim)
``H2 holds in full'' is false: the combination is NEGATIVE for the core test, NOISY for collusion, and NULL for cold start.
% src: A9 (verbatim); F17

\begin{figure}[t]
\centering
\includegraphics[width=\columnwidth]{\figdir/fig5_s4_gap_vs_correlation.pdf}
\caption{Diagnostic AUC of D minus C against the nominal correlation $\rho$ (a) and the realised pair concentration (b), in this simulation; 5 seeds, 95\% intervals over seeds. The pre-registered floor of $+0.02$ is the line a gain of D over C would have to reach. (Generated by \texttt{scripts/make\_figures.py}.)}
\label{fig:gap}
\end{figure}

\paragraph*{The closed loop agrees}
In this simulation, D is not better than C on any closed-loop point: the paired success difference D minus C is $-0.016$ [$-0.027$, $-0.006$] at the shared point (NEGATIVE at 20\% and 30\%, NULL at 5\% and 10\%), and it is about the same at $\rho=0$ and $\rho=1.0$ ($-0.016$), so it does not widen with $\rho$.
% src: A10 (F26, F27 item 4; D56 section 3)
This is the secondary closed-loop check of H2 and does not revise the shared-onset verdict above.
Under staggered onset (below) the direction is the same.
% src: A10 (F40, descriptive)

\paragraph*{Why, as far as it was measured}
The \texttt{same\_segment} flag is correct on every sampled RSU--RSU edge (0 mismatches over 92,000 instances, and 0 on vehicle--RSU edges), and D's output depends on it (mean $|\Delta\text{trust}|$ of $0.024$--$0.103$ when it is zeroed, set to 1 or permuted), but removing it does not systematically cost detection.
% src: A11 (F18, D48)
A degraded RSU's own evidence moves at the same time as its degraded same-segment neighbours' (own-lead $-0.3$ to $+0.5$\,s, all intervals include $0$; about 110\,s after onset), because onset is shared.
At $\rho=0$, 59\% of degraded RSUs have only healthy same-segment neighbours; at $\rho=1$, 95\% have degraded neighbours, but those move simultaneously.
% src: A11 (F18, F19)
So in this scenario own-node evidence is never later than neighbour evidence, and a lead-time advantage for message passing is not available by construction.
We state this as a \emph{boundary condition of this scenario}, not as a general result.
% src: A11 (verbatim)

\paragraph*{Staggered onset (S13, Experiment 2)}
To create a regime in which neighbour evidence \emph{does} lead, each degraded RSU's onset is drawn independently and uniformly over $[400,700]$ ($S=300$).
In this simulation, a late-onset degraded RSU (one with an earlier-onset degraded same-segment RSU--RSU neighbour) showed degraded evidence in \texttt{success\_ewma} after its earlier-onset neighbours by $+92.3$ [$+79.2$, $+105.5$] steps at $\rho=1.0$ and $+103.6$ [$+70.9$, $+136.2$] at $\rho=0$ (draws 50--79, seeds 1--5; mean onset gap $98.1$ and $103.6$).
% src: A24 (F35, D65 2.3)
The lead is the injected stagger carried through the observable; it is a statement about the scenario, not about any model.
With fresh heads trained on the staggered regime (draws 200--249), D detected late-onset RSUs \emph{later} than C: D minus C is $+70.5$ [$+42.8$, $+98.2$] steps at $\rho=1.0$ (C $126.5$, D $197.0$) and $+96.3$ [$+66.8$, $+125.9$] at $\rho=0$ (C $130.1$, D $226.4$).
% src: A25 (F36-F38, F40; D65, D66)
The pre-registered label is NEGATIVE at $\rho=1.0$ (C beats D) and the $\rho=0$ control is NEGATIVE too, so there is no structural reading; the extension to seeds 6--10 was not triggered.
Descriptively, the window AUC difference D minus C is $-0.046$ [$-0.100$, $+0.008$] at $\rho=1.0$ and $-0.127$ [$-0.204$, $-0.050$] at $\rho=0$, and closed-loop post-onset success D minus C is $-0.016$ [$-0.021$, $-0.011$] at $\rho=1.0$ and $-0.017$ [$-0.027$, $-0.008$] at $\rho=0$, each model with its own weights tuned on seed 100.
% src: A25 (descriptive, no pre-registered label)
D42 accounting: at $\rho=1.0$, 250 draws attempted, 0 failed, 250 aggregated; at $\rho=0$, 250 / 0 / 164 (draws with no late-onset RSU are not aggregated).
Latency is compared under one shared rule, not at a detection operating point; the rule's false-alarm rate was 0.32--0.44 of healthy RSUs where measured (S16) and was not measured elsewhere.
% src: A25 (D78 shared-rule caveat, verbatim)
The graph did not help in this regime either, where a neighbour lead exists.
This does not revise the shared-onset result; it is a separate, labelled regime.
% src: A25 (verbatim sentence); X17 (no generalisation); X19 (D's smaller deficit at rho 1.0 is not reported as an advantage)

\paragraph*{A contained graph head (S14)}
The S14 experiment tested one candidate explanation, as an exploratory, pre-registered variant: D2 is C's node-local branch plus a mean over same-segment RSU neighbours only (no vehicle nodes), with neighbour weights initialised to zero (1570 parameters against C's 1313).
In this simulation, D2 was not distinguishable from C.
% src: A28; D68
The D2 minus C detection latency of late-onset RSUs under staggered onset is $-1.7$ [$-7.0$, $+3.5$] steps at $\rho=1.0$ and $-0.8$ [$-4.0$, $+2.5$] at $\rho=0$, both labelled NULL, so a $20$-step difference in either direction is ruled out.
The shared-onset diagnostic AUC difference is $+0.0012$ [$-0.0018$, $+0.0043$] at $\rho=1.0$ and $-0.0008$ [$-0.0043$, $+0.0026$] at $\rho=0$ (NULL), and the closed-loop post-onset success difference is $+0.0003$ [$-0.0002$, $+0.0009$] in the shared regime and $+0.0004$ [$-0.0001$, $+0.0008$] in the staggered regime (NULL), all with 95\% intervals over 5 seeds (draws 250--299).
% src: A28 (F42-F45)
Table~\ref{tab:paper_d2_dispatch} gives \nav, \nwin\ and the never-avoided rate together for A, C and D2 in both regimes.
Draws attempted were 250 and failed 0 at every point, and 250 were aggregated except for the $\rho=0$ latency control (183, since draws with no late-onset RSU are not aggregated).
% src: A28 (D42)
Latency is compared under one shared rule, not at a detection operating point; the rule's false-alarm rate was 0.32--0.44 of healthy RSUs where measured (S16) and was not measured elsewhere.
% src: A28 (D78 shared-rule caveat, verbatim)
Descriptively, the unconstrained D trailed C under staggered onset (draws 200--249), while D2, which keeps C's own-node branch, matched C (draws 250--299); D and D2 were not compared directly on the same draws.
% src: A28 (last sentences); X21, X22 (the opposite-signed n_av differences between C and D2 are not reported as an effect)

\input{\tabdir/paper_d2_dispatch.tex}

\subsection{Training objective}
\label{sec:objective}

Every head above is trained on the next-task objective.
We departed from it once, in a pre-registered exploratory test (D72): C-h and D2-h are C's and D2's architectures trained on a forward-looking objective, with label 1 if the fraction of an RSU's tasks meeting their deadline in the next 150 steps is below $0.75$, and windows with fewer than 5 tasks excluded.
A pre-registered confirmatory test (D76) of one descriptive observation from that experiment followed.
Table~\ref{tab:paper_objective} gives both, with their labels.
% src: A29, A30; D72, D76; R6, R7

\input{\tabdir/paper_objective.tex}

\paragraph*{A graph head under the horizon objective (S15)}
In this simulation, under staggered onset ($S=300$; draws 300--349, 5 seeds), D2-h was not distinguishable from C-h: D2-h minus C-h detection latency of late-onset RSUs is $-4.1$ [$-12.0$, $+3.7$] steps at $\rho=1.0$ and $-2.0$ [$-8.1$, $+4.1$] at $\rho=0$, both labelled NULL, so a $20$-step difference in either direction is ruled out.
% src: A29 (F48, F49; D72)
D42 accounting: at $\rho=1.0$, 250 draws attempted, 0 failed, 250 aggregated; at $\rho=0$, 250 / 0 / 166.
Descriptively, the closed-loop post-onset success difference D2-h minus C-h is $+0.0020$ [$+0.0007$, $+0.0033$] (no label), with \nav\ / \nwin\ / never-avoided $589.6$ / $903.7$ / $0.468$ for A, $224.8$ / $559.9$ / $0.020$ for C-h and $224.0$ / $556.9$ / $0.020$ for D2-h.
% src: A29 (D42; descriptive closed loop; R2: always together); F48
D2-h's output depended on its neighbour block in 4 of 5 evaluation seeds (seed 1 did not: $0.0014$).
Latency is compared under one shared rule, not at a detection operating point; the rule's false-alarm rate was 0.32--0.44 of healthy RSUs where measured (S16) and was not measured elsewhere.
% src: A29 (F46 probe; D78 shared-rule caveat, verbatim); X23 (no claim about other horizons, thresholds or architectures)

\paragraph*{The horizon objective and the node-local head's detection (S16)}
In S15, descriptively, C-h detected late-onset RSUs earlier than the next-task head C (draws 300--349: $37.4$ steps earlier [$21.2$, $53.5$] at $\rho=1.0$ and $34.9$ [$14.1$, $55.8$] at $\rho=0$); that observation stands as recorded.
% src: A30 (S15 descriptive observation); F48, F49
In this simulation, in a pre-registered confirmatory test (D76; fresh draws 350--399, staggered onset $S=300$, fraction $0.20$, 5 seeds, all degraded RSUs, each model its own pre-onset 1st-percentile threshold), C-h detected degraded RSUs $38.2$ steps earlier than C [$23.1$, $53.4$] at $\rho=1.0$ and $35.7$ steps earlier [$18.1$, $53.4$] at $\rho=0$ (EARLIER at both).
The C-h minus C false-alarm rate on healthy RSUs was $-0.127$ [$-0.460$, $+0.207$] at $\rho=1.0$ and $-0.117$ [$-0.467$, $+0.232$] at $\rho=0$ (NOT ACCEPTABLE at both: more than 5 percentage points more false alarms is not ruled out), so the pre-registered verdict is NOT CONFIRMED.
% src: A30 (F50, F51; D76); X24 (EARLIER only beside NOT ACCEPTABLE and NOT CONFIRMED)
False-alarm rates under the rule are $0.316$ for C-h and $0.443$ for C ($\rho=1.0$).
Descriptively, closed-loop success C-h minus C is $-0.0008$ [$-0.0023$, $+0.0007$], with \nav\ / \nwin\ / never-avoided $589.5$ / $907.1$ / $0.440$ for A, $245.4$ / $576.5$ / $0.024$ for C and $237.2$ / $565.0$ / $0.024$ for C-h.
D42 accounting: 250 draws attempted, 0 failed, 250 aggregated at every point.
We offer no mechanism.
Latency is compared under one shared rule, not at a detection operating point; the rule's false-alarm rate was 0.32--0.44 of healthy RSUs where measured (S16) and was not measured elsewhere.
% src: A30 (rates, closed loop, R2 triplet, D42, "No mechanism"; D78 shared-rule caveat, verbatim); X23, X24

\subsection{Explanation fidelity and overhead}
\label{sec:overhead}

\paragraph*{Fidelity (H4)}
Table~\ref{tab:h4} reports both passes of the pre-registered Layer 2 test, each with its own literal label.
In this simulation, the first pass (draw 100) gives a trust-change gap of $+0.160$ [$+0.043$, $+0.277$] (WEAK) and a decision-change gap of $+0.033$ [$+0.0003$, $+0.065$] (WEAK); the replication (draw 101) gives $+0.158$ [$+0.031$, $+0.284$] (WEAK) and $+0.0500$ [$+0.017$, $+0.083$], which is FAITHFUL by the literal $\ge$ rule \emph{exactly at the $0.05$ floor} and which we therefore call borderline.
Top and random-low rates accompany every gap in Table~\ref{tab:h4}.
% src: A13, R3; F24, F25, F26/F27 (H4 block)
We make no stronger statement than ``weakly faithful'' for Layer 2: the masks are unstable (agreement $0.22$ against $0.80$), so each explanation is one draw with its run seed, and the neighbour clause is conditional and rare (10\% of D's explanations; 16\% when the explained RSU is degraded).
% src: A13, X5

\paragraph*{Overhead (H5)}
Table~\ref{tab:paper_latency} gives the median time per decision by component, measured separately, on a 13th Gen Intel i5-13450HX (CPU only, one decision at a time).
% src: A14; results/s10/tables/paper_latency, latency_breakdown
At one thread the median is $0.27$\,ms for trust inference, $0.04$\,ms for selection, $0.09$\,ms for Layer 1 and $310$\,ms for Layer 2; at ten threads it is $0.94$, $0.05$, $0.11$ and $675$\,ms.
Layer 2 is faster on one thread for graphs this small.
Against the stated $100$\,ms assumption, Layer 2 is asynchronous only.
% src: A14 (numbers verbatim), A13 ("asynchronous only")

\input{\tabdir/paper_latency.tex}
